% !TEX root = main.tex
% ---------------------------------------------------------------------------
% Three paper outlines. The Method section (method.tex) is SHARED by all three;
% what changes is which subsection carries the contribution, and which
% experiments are promoted to the headline table.
% Status markers:  [HAVE] measured   [RUN] in flight   [NEED] not yet run
% ---------------------------------------------------------------------------
\section{Paper outlines}
\label{sec:outlines}

All three versions share \S\ref{sec:method} and differ in framing, headline experiment, and
which related-work axis they engage. A shared claim ledger is given in
\S\ref{sec:outlines:ledger}; per-version deltas are noted inline.

% ===========================================================================
\subsection{Version A --- ``Coverage, not more demonstrations''}
\label{sec:outlines:A}

\paragraph{Titles.}
\emph{Coverage over Count: Autonomous Stress-Aware Demonstrations for Gentle Robotic
Grasping}\,/\,\emph{More Demonstrations Do Not Fix Precise Grasping}.

\paragraph{Abstract sketch.}
Gentle grasping of fragile food is a precision problem: the correct gripper width is a
function of pose, size, shape and stiffness, none directly observable. We show that scaling
teleoperated-style demonstrations does not resolve it---success saturates while the failure
mode remains incorrect grasp width---and that what is scarce is coverage of the physical
variation rather than trajectory count. We therefore generate demonstrations autonomously
with a stress-minimising grasp planner over a randomised scene distribution, and show
on hardware that broad simulated coverage co-trained with a small real slice outperforms both
pure-simulation and real-only policies, while simulated ranking alone fails to predict real
ranking.

\paragraph{Section outline.}
\begin{enumerate}
  \item \textbf{Introduction.} Precision framing; the saturation observation as the hook;
        contributions (i) the coverage diagnosis, (ii) the stress-aware autonomous
        demonstrator, (iii) the transfer recipe and its real validation.
  \item \textbf{Related work.} Demonstration scaling and data-quality work; sim2real for
        deformables; grasp-quality metrics (force-closure vs.\ stress-based); gentle/compliant
        manipulation (tactile-centric) as the contrasting tradition.
  \item \textbf{Motivating study: does demonstration count help?} \emph{Own section, not a
        result buried in ablations.} Fixed architecture, three dataset sizes, failure-mode
        decomposition (approach vs.\ width). Extend to real demonstration counts
        \textbf{[RUN]}.
  \item \textbf{Method} (\S\ref{sec:method}), emphasis on
        \S\ref{sec:method:synth}--\ref{sec:method:demos}: the demonstrator is the coverage
        instrument.
  \item \textbf{Experiments.}
    \begin{enumerate}\renewcommand{\labelenumii}{\alph{enumii})}
      \item Demonstrator quality: planning success, reproducibility, stress vs.\ a geometric
            baseline \textbf{[HAVE]}.
      \item Simulated policy performance and the demonstrator-property ablations (firmness,
            occlusion bound, approach profile, stop supervision) \textbf{[HAVE/RUN]}.
      \item Coverage ablations: size band, shape/mesh pool, out-of-distribution size
            \textbf{[RUN]}.
      \item Real-robot: pure-sim vs.\ real-only vs.\ co-trained; real-demo-count curve
            \textbf{[HAVE/RUN]}.
      \item Generalist: single policy across categories \textbf{[RUN]}.
    \end{enumerate}
  \item \textbf{Analysis.} Why simulated ranking inverts on hardware; the measured
        perception offset; what remains unexplained.
  \item \textbf{Limitations and conclusion.}
\end{enumerate}

\paragraph{Key figures/tables.}
(F1) saturation curve + failure-mode split; (F2) pipeline schematic; (F3) demonstrator stress
distributions vs.\ baseline; (T1) simulated ablation table; (T2) real-robot table with
sample sizes; (F4) real-demo-count curve; (F5) qualitative real rollouts.

\paragraph{Version-specific risk.} A reviewer reads the saturation result as
architecture-specific. Mitigation: pair it with the demonstrator-property ablation, where the
\emph{same} architecture and count move substantially once the demonstrations change---which
is precisely the claim that count is not the binding constraint.

% ===========================================================================
\subsection{Version B --- ``Gentleness without touch''}
\label{sec:outlines:B}

\paragraph{Titles.}
\emph{Gentleness Without Touch: Compiling Simulated Stress into Tactile-Free Manipulation
Policies}\,/\,\emph{Stress as Supervision}.

\paragraph{Abstract sketch.}
Delicate manipulation is usually built around tactile sensing, the closest observable proxy
for the internal stress that actually damages an object. We observe that this quantity is not
merely observable but exact inside a deformable simulator, and use it there: a
stress-minimising planner selects grasps by predicted internal stress under a
position-controlled contact model, and the resulting demonstrations are distilled into a
policy that acts from a single depth view and proprioception, with no contact sensing at
deployment. We evaluate on a real arm across object poses, sizes and shapes, and compare
against a tactile-equipped policy on the same task.

\paragraph{Section outline.}
\begin{enumerate}
  \item \textbf{Introduction.} Touch as the default modality and its practical costs; the
        observation that the causal quantity is exact in simulation; the claim that
        gentleness can be \emph{compiled into demonstrations} rather than sensed.
  \item \textbf{Related work.} Tactile manipulation and tactile-conditioned policies;
        privileged-information / teacher--student distillation; deformable simulation;
        stress- and fragility-aware grasp metrics.
  \item \textbf{Method} (\S\ref{sec:method}), emphasis on \S\ref{sec:method:synth}: the
        contact model, the $E$-linearity that makes material randomisation cheap, and the
        score terms as an operational definition of ``gentle''.
  \item \textbf{Experiments.}
    \begin{enumerate}\renewcommand{\labelenumii}{\alph{enumii})}
      \item Does the metric mean anything? Stress-minimising vs.\ geometric planner, at
            matched success \textbf{[HAVE]}; sensitivity to material and size.
      \item Does gentleness survive distillation? Policy-induced stress vs.\ demonstrator
            stress in simulation \textbf{[HAVE, needs a dedicated presentation]}.
      \item Real-robot task performance without touch \textbf{[HAVE]}, across pose/size/shape
            and categories \textbf{[RUN]}.
      \item \textbf{Tactile comparison} \textbf{[NEED]}: a tactile-conditioned policy trained
            on real demonstrations on the same task and protocol. \emph{This is the
            experiment that makes the framing defensible.}
      \item Real gentleness proxy \textbf{[NEED]}: post-lift deformation or contact-force
            logging, so ``gentle'' is not exclusively a simulated claim.
    \end{enumerate}
  \item \textbf{Analysis.} What the stress signal buys (which failure modes it removes:
        stem/rim grasps, over-squeeze); where it is blind (dynamic effects during lift,
        which the quasi-static model does not represent).
  \item \textbf{Limitations and conclusion.} Explicitly: we do not claim touch is unnecessary
        in general.
\end{enumerate}

\paragraph{Key figures/tables.}
(F1) the idea figure---stress field inside the simulated object next to the real camera view,
same grasp; (F2) planner comparison with stress distributions and the bruising threshold
marked; (F3) pipeline; (T1) real results, tactile-free vs.\ tactile baseline; (F4)
grasp-selection gallery showing rejected stem/rim candidates and the accepted envelop; (T2)
gentleness metrics, simulated and real proxy.

\paragraph{Version-specific risk.} Without (d) the paper asserts a tactile-free advantage it
never tests. Two pre-registered answers to the objections already in our notes: \emph{if
tactile is unavailable, is depth unavailable too?}---the argument is cost, fingertip geometry,
wear and per-gripper recalibration, not sensor availability in principle; \emph{why not
simulate tactile instead?}---because tactile is a proxy for stress, and in simulation we can
supervise the causal quantity directly.

% ===========================================================================
\subsection{Version C --- ``What it takes to transfer a precise, gentle skill''}
\label{sec:outlines:C}

\paragraph{Titles.}
\emph{Fixed Points and Blind Spots: An Empirical Account of Sim-to-Real Transfer for Precise
Gentle Manipulation}\,/\,\emph{What It Takes}.

\paragraph{Abstract sketch.}
We report an end-to-end study of transferring a precise, gentle grasping skill from
deformable simulation to a real arm. Beyond the working recipe, we isolate two failure modes
of behaviour cloning with absolute action chunks---supervision without lead, and trajectory
dwell---each of which yields excellent training loss and near-zero success, and give dataset
predicates that detect both before training. We further show that simulated ranking does not
transfer across data regimes, and quantify the perception component of the sim-to-real gap
using step-aligned real/simulated observation pairs.

\paragraph{Section outline.}
\begin{enumerate}
  \item \textbf{Introduction.} The narrow-path argument; contributions as \emph{findings}
        rather than a system.
  \item \textbf{System} (condensed \S\ref{sec:method}) --- enough to make the findings
        reproducible.
  \item \textbf{Finding 1: absolute-action BC has two fixed points.}
        \S\ref{sec:method:supervision} promoted to a full section: mechanism, minimal
        reproduction, the $2\times2$ control (representation $\times$ derivation) that
        exonerates every other factor, the two predicates, and the recovery
        ($0 \rightarrow$ working) when the dwell condition is met \textbf{[HAVE]}.
  \item \textbf{Finding 2: demonstration properties dominate demonstration count.}
        Saturation study plus the firmness/approach/stop ablations \textbf{[HAVE/RUN]}.
  \item \textbf{Finding 3: simulated ranking inverts on hardware.} The three-way real
        comparison and what it implies for model selection \textbf{[HAVE]}.
  \item \textbf{Finding 4: the gap is partly a measurable perception offset.} The paired
        real/simulated replay methodology---which is itself a reusable contribution---and the
        residual decomposition \textbf{[HAVE]}.
  \item \textbf{Putting it together.} The final recipe; specialist and generalist results
        \textbf{[RUN]}.
  \item \textbf{Limitations and conclusion.}
\end{enumerate}

\paragraph{Key figures/tables.}
(F1) the stall: training loss vs.\ evaluation success for the failing and fixed datasets,
with the action-difference histogram inset; (F2) the $2\times2$ control table; (F3) paired
real/simulated cloud overlay before and after the offset correction; (T1) the recipe table
(component $\rightarrow$ measured effect); (T2) real-robot results.

\paragraph{Version-specific risk.} Perceived as an experience report. Mitigation: Finding~1 is
architecture-agnostic and, to our knowledge, undocumented; lead with a minimal reproduction
that a reader could run on their own BC stack, and state the predicates as reusable dataset
gates.

% ===========================================================================
\subsection{Shared claim ledger}
\label{sec:outlines:ledger}

Every claim any of the three versions would make, with the evidence backing it and its
current status.

\begin{center}\small
\begin{tabular}{p{0.46\linewidth} p{0.30\linewidth} l}
\toprule
\textbf{Claim} & \textbf{Evidence} & \textbf{Status} \\
\midrule
Demonstration count saturates; failure mode is grasp width & three dataset sizes, fixed
architecture, simulated & \textbf{[HAVE]} (sim only) \\
Real-demonstration count curve & $N\!\in\!\{0,5,10,50\}$ co-training, real trials & \textbf{[RUN]} \\
Stress-minimising planner is gentler than a geometric one & paired collections, matched
success & \textbf{[HAVE]} \\
Demonstrator firmness moves policy success & controlled recipe change & \textbf{[HAVE]} \\
Occlusion bound reduces hidden-object observations & ground-truth occlusion measure & \textbf{[HAVE]} \\
Approach-profile / stop-frame effects & v3.1--v3.3 recipe series & \textbf{[RUN]} \\
Absolute-action BC needs lead; dwell is a second, independent stall & controlled datasets,
$2\times2$ & \textbf{[HAVE]} \\
Simulated ranking inverts on hardware & three policies, real trials & \textbf{[HAVE]}, small $n$ \\
Perception offset measured and correctable & step-aligned real/sim pairs & \textbf{[HAVE]} \\
Co-training beats pure-sim and real-only in real & real trials & \textbf{[HAVE]}, small $n$ \\
Generalist across categories & pooled multi-category training & \textbf{[RUN]} \\
Policy is gentle \emph{on real produce} & --- & \textbf{[NEED]} \\
Tactile-free is competitive with tactile & --- & \textbf{[NEED]} (Version~B only) \\
\bottomrule
\end{tabular}
\end{center}

\paragraph{Recommendation.} Version~A is publishable with what exists plus the in-flight runs.
Version~B is the most distinctive but is gated on one new experiment (the tactile baseline)
and one measurement (a real gentleness proxy); if both are run it is the strongest paper.
Version~C is the safest technically and should in any case supply the experiment sections of
A or B---Findings~1 and~4 are the most transferable results we have. The framing to avoid in
all three is a blanket ``gentle manipulation of real food'' claim: gentleness is currently
validated in simulation, task success on hardware, and the link between them is the missing
measurement.
